\section{Detailed benchmark results and reaction-path analysis}
\label{sec:si-detailed-results}

\subsection{Complete workflow and tier comparisons}

The dataset contains 200 task runs: one per task under Guided, Standard, Record, and Brief. Each run has a composite and all three component scores. Table~\ref{tab:si-arm-tier} reports overall and tier means. Standard uses the same model backbone as Guided; Record removes required physical validation, and Brief shortens diagnostic feedback.

\begin{table}[!htbp]
\centering
\small
\caption{Composite and component means for each workflow and task tier.}
\label{tab:si-arm-tier}
\begin{tabular}{llrrrrr}
\toprule
Tier & Workflow & $n$ & Composite & Graph & SC & Accuracy \\
\midrule
All & Guided & 50 & 0.7930 & 0.8851 & 0.9233 & 0.7124 \\
 & Standard & 50 & 0.5814 & 0.7774 & 0.5303 & 0.5743 \\
 & Record & 50 & 0.6918 & 0.7813 & 0.7177 & 0.6639 \\
 & Brief & 50 & 0.7541 & 0.8493 & 0.8468 & 0.6920 \\
\addlinespace
Easy & Guided & 20 & 0.9125 & 0.9778 & 0.9250 & 0.8954 \\
 & Standard & 20 & 0.6989 & 0.8849 & 0.5858 & 0.7244 \\
 & Record & 20 & 0.9128 & 0.9419 & 0.9050 & 0.9118 \\
 & Brief & 20 & 0.9370 & 0.9681 & 0.9125 & 0.9440 \\
\addlinespace
Normal & Guided & 10 & 0.7370 & 0.8589 & 0.9021 & 0.6342 \\
 & Standard & 10 & 0.5731 & 0.7607 & 0.5217 & 0.5675 \\
 & Record & 10 & 0.5287 & 0.6786 & 0.5783 & 0.4789 \\
 & Brief & 10 & 0.5857 & 0.7472 & 0.7740 & 0.4646 \\
\addlinespace
Hard & Guided & 10 & 0.9577 & 0.9107 & 0.9542 & 0.9672 \\
 & Standard & 10 & 0.6887 & 0.8481 & 0.5033 & 0.7549 \\
 & Record & 10 & 0.8723 & 0.9243 & 0.6650 & 0.9672 \\
 & Brief & 10 & 0.8282 & 0.9286 & 0.7767 & 0.8372 \\
\addlinespace
Impossible & Guided & 10 & 0.4451 & 0.7006 & 0.9100 & 0.1700 \\
 & Standard & 10 & 0.2473 & 0.5083 & 0.4550 & 0.1000 \\
 & Record & 10 & 0.2325 & 0.4200 & 0.5350 & 0.0500 \\
 & Brief & 10 & 0.4829 & 0.6345 & 0.8583 & 0.2700 \\
\bottomrule
\end{tabular}
\par\smallskip\begin{minipage}{\linewidth}\footnotesize SC denotes the graded self-consistency component. Every mean is calculated from the complete publication table. The overall tier weights are 20/50 for Easy and 10/50 for each other tier.\end{minipage}
\end{table}


Table~\ref{tab:si-pairwise} gives all six pairwise comparisons on the same 50 tasks. Guided--Standard has a mean difference of 0.211564, a median difference of 0.114200, and 36 task wins, seven losses, and seven ties. The 95\% task-bootstrap interval is [0.128988, 0.298537]. Within the Normal tier, the means are 0.737020 and 0.573050, with a difference of 0.163970 [0.036520, 0.292230] and seven wins in ten tasks. Within Hard, the means are 0.957650 and 0.688750; the difference is 0.268900 [0.121300, 0.445930], with Guided higher in all ten task pairs. Easy has means of 0.912520 and 0.698910, with 12 wins, four losses, and four ties; its paired median difference is 0.050000. Impossible has means of 0.445060 and 0.247330 and is examined at the level of the requested target below.

\begin{table}[!htbp]
\centering
\small
\caption{All pairwise composite-score comparisons across the 50 tasks.}
\label{tab:si-pairwise}
\begin{tabular}{lrrr}
\toprule
Comparison & Mean difference & 95\% interval & W/L/T \\
\midrule
Guided--Record & +0.1012 & [+0.0396, +0.1695] & 28/7/15 \\
Guided--Standard & +0.2116 & [+0.1290, +0.2985] & 36/7/7 \\
Guided--Brief & +0.0388 & [-0.0317, +0.1076] & 26/10/14 \\
Record--Standard & +0.1104 & [+0.0167, +0.2062] & 28/13/9 \\
Record--Brief & -0.0623 & [-0.1271, +0.0001] & 14/24/12 \\
Standard--Brief & -0.1727 & [-0.2713, -0.0762] & 11/34/5 \\
\bottomrule
\end{tabular}
\par\smallskip\begin{minipage}{\linewidth}\footnotesize The difference is the first workflow minus the second. W/L/T gives the first workflow's task wins, losses, and ties. Percentile task-bootstrap intervals use 100,000 resamples and seed 20261031; the four configurations of each sampled task remain together.\end{minipage}
\end{table}


Table~\ref{tab:si-weighted-gain} decomposes the Guided--Standard difference into weighted score components. Across all tasks, self-consistency contributes +0.117875 and accuracy +0.082916. In Normal and Hard together, the composite difference is +0.216435, including +0.124688 from self-consistency and +0.083712 from accuracy. In Hard, accuracy contributes +0.127391 and self-consistency +0.135250: both the requested answers and the calculations supporting them contribute to the gain.

\begin{table}[!htbp]
\centering
\small
\caption{Weighted decomposition of the Guided--Standard composite difference.}
\label{tab:si-weighted-gain}
\begin{tabular}{lrrrrr}
\toprule
Scope & $n$ & Composite & $0.1\Delta G$ & $0.3\Delta SC$ & $0.6\Delta A$ \\
\midrule
All & 50 & +0.211564 & +0.010776 & +0.117875 & +0.082916 \\
Easy & 20 & +0.213610 & +0.009287 & +0.101750 & +0.102579 \\
Normal & 10 & +0.163970 & +0.009823 & +0.114125 & +0.040033 \\
Hard & 10 & +0.268900 & +0.006259 & +0.135250 & +0.127391 \\
Impossible & 10 & +0.197730 & +0.019226 & +0.136500 & +0.042000 \\
Normal + Hard & 20 & +0.216435 & +0.008041 & +0.124688 & +0.083712 \\
\bottomrule
\end{tabular}
\par\smallskip\begin{minipage}{\linewidth}\footnotesize Component contributions are calculated from the unrounded final components. Composite differences use the four-decimal final scores, so the last displayed digits can differ by score rounding.\end{minipage}
\end{table}


Guided--Record averages +0.101156 [0.039606, 0.169478] across all tasks. In Normal and Hard together, its +0.146845 composite difference decomposes into +0.008338 from graph, +0.091938 from self-consistency, and +0.046574 from accuracy. Guided--Brief averages +0.038816 [$-0.031704$, 0.107648], with mean differences of +0.151350 on Normal and +0.129470 on Hard. Individual multistep tasks differ: NB1 has composite scores of 0.9628, 0.9128, and 0.9900 for Guided, Standard, and Brief; Guided and Brief both have self-consistency 1. NI1 has Guided and Standard composites of 0.9958 and 0.6196, with self-consistency 1 and 0.35. NQ5 favors Record over Guided, while NU3 favors Brief. The leading configuration varies among multistep tasks.

\subsection{Graph delivery and evaluation of the run results}

During execution, Guided, Record, and Brief each wrote 50 root mechanism graphs; Standard wrote 43. Standard's remaining seven studies---F1, HA1, HF1, IA9, M17, NR1, and P3---were organized from their own conversations, reports, tool results, and calculation files. The common evaluation procedure thus produced 50 graphs per workflow and 200 complete score profiles, with the source of each recovered field recorded. The recovered graphs give Standard credit for the research documented in these runs. This reduces the Guided--Standard difference compared with scoring by final-file delivery alone.

The seven Standard studies have a mean organized composite of 0.153743, compared with 0.828529 for Guided on the same tasks. They account for 44.7\% of the total Guided--Standard difference. Table~\ref{tab:si-sc-distribution} separately reports self-consistency for the 40 non-Impossible tasks. Guided has 24 full-score records (60.0\%) and 16 partial-score records; Standard has seven full-score records (17.5\%), 29 partial-score records, and four zero-score records.

\begin{table}[!htbp]
\centering
\small
\caption{Self-consistency distribution over all 40 non-Impossible tasks.}
\label{tab:si-sc-distribution}
\begin{tabular}{lrrrr}
\toprule
Workflow & SC $=1$ & $0<$ SC $<1$ & SC $=0$ & Mean SC \\
\midrule
Guided & 24 & 16 & 0 & 0.9266 \\
Standard & 7 & 29 & 4 & 0.5492 \\
Record & 16 & 24 & 0 & 0.7633 \\
Brief & 20 & 20 & 0 & 0.8439 \\
\bottomrule
\end{tabular}
\par\smallskip\begin{minipage}{\linewidth}\footnotesize This table describes the graded component for the positive-chemistry tiers. The separate scientific-validation record is evaluated from the complete task contract.\end{minipage}
\end{table}


\subsection{Elapsed research time and accumulated calculation effort}

Table~\ref{tab:si-batch-times} preserves every batch observation from the 48-run repeated-task experiment. A batch contains the same four tasks at a concurrency of two, and its elapsed time spans the batch's actual start and finish. Guided's mean is 4.4514 h, Standard's 3.9577 h, Record's 3.8942 h, and Brief's 4.8273 h. Relative to Standard, Guided's observed mean is 12.5\% longer; relative to Brief, it is 7.8\% shorter. These observations include waiting, calculation execution, and continuation of the research after results return.

\begin{table}[!htbp]
\centering
\small
\caption{All 12 elapsed-time observations in the repeated-task study, in hours.}
\label{tab:si-batch-times}
\begin{tabular}{lrrrr}
\toprule
Workflow & Batch 1 & Batch 2 & Batch 3 & Mean \\
\midrule
Guided & 4.2383 & 3.8669 & 5.2489 & 4.4514 \\
Standard & 5.3381 & 2.5044 & 4.0306 & 3.9577 \\
Record & 5.4189 & 2.6203 & 3.6433 & 3.8942 \\
Brief & 7.2689 & 3.6475 & 3.5656 & 4.8273 \\
\bottomrule
\end{tabular}
\par\smallskip\begin{minipage}{\linewidth}\footnotesize Each observation is a four-task batch (E1, NB1, HA1, IA1) at concurrency two. Three batches per workflow contain 48 task runs in total. Batch completion times include the original research lifecycle and outcomes.\end{minipage}
\end{table}


The P2 effort comparison uses the common 15-task set D3, M1, M8, NB1, NQ5, NR1, NU3, NU5, NU6, HA1, HA6, HB1, HB10, HD2, and HF1. This yields 60 run-level effort records, with all four configurations represented for each task. Table~\ref{tab:si-common-effort} gives the total and mean job-hours, job counts, terminal outcomes, and corresponding final scores. Recorded ledger durations are summed across quantum-chemical jobs. For Standard/M1, the archived effort reconstruction uses native-log CPU time divided by four, as specified in the source record. The measurement unit throughout this comparison is accumulated job-hours.

\begin{table}[!htbp]
\centering
\small
\caption{Recorded quantum-chemical effort on the 15 tasks shared by all configurations.}
\label{tab:si-common-effort}
\begin{tabular}{lrrrrrr}
\toprule
Workflow & Job-h total & Job-h/task & Jobs & Failed & Cancelled & Mean $S$ \\
\midrule
Guided & 166.19 & 11.079 & 789 & 114 & 94 & 0.8124 \\
Standard & 113.32 & 7.555 & 419 & 45 & 50 & 0.6495 \\
Record & 118.16 & 7.877 & 431 & 62 & 38 & 0.7360 \\
Brief & 233.85 & 15.590 & 883 & 167 & 120 & 0.7273 \\
\bottomrule
\end{tabular}
\par\smallskip\begin{minipage}{\linewidth}\footnotesize Each workflow contributes 15 task runs. Job-hours sum recorded process durations. Failed and cancelled jobs are separate terminal categories. Job counts measure execution effort; scientific validation is assessed from the calculation results.\end{minipage}
\end{table}


Guided used a mean of 11.0793 job-hours and obtained a mean composite of 0.812413. Relative to Standard, it used 46.7\% more job time and scored 0.162933 higher; relative to Brief, it used 28.9\% less job time and scored 0.085080 higher. Guided had both a higher score and lower job time than Brief on seven of the 15 paired tasks. Guided recorded 789 jobs, including 114 failures and 94 cancellations; Brief recorded 883 jobs, including 167 failures and 120 cancellations. In the wider 63-run archive, NU3, NU5, NU6, and NR1 account for 163.5 of 673.82 job-hours across their 16 runs. These trajectories record repeated searches for the required intermediates and connecting pathways.

\subsection{Conformer references and reaction-step assignments}

Table~\ref{tab:si-case-profiles} gives all score components for the following cases. We examine which reactant conformer defines each barrier and which reaction step each optimized saddle connects.

\begin{table}[!htbp]
\centering
\small
\caption{Complete score profiles for the four scientific-attribution cases.}
\label{tab:si-case-profiles}
\begin{tabular}{llrrrr}
\toprule
Task & Workflow & Composite & Graph & SC & Accuracy \\
\midrule
HA8 & Guided & 0.9625 & 1.0000 & 0.8750 & 1.0000 \\
 & Standard & 0.6250 & 1.0000 & 0.3500 & 0.7000 \\
HH1 & Guided & 0.9838 & 1.0000 & 1.0000 & 0.9731 \\
 & Standard & 0.8338 & 1.0000 & 0.5000 & 0.9731 \\
HI1 & Guided & 0.9999 & 0.9992 & 1.0000 & 1.0000 \\
 & Standard & 0.7032 & 0.3325 & 0.2333 & 1.0000 \\
NU6 & Guided & 0.6719 & 1.0000 & 0.9375 & 0.4844 \\
 & Standard & 0.6892 & 0.8167 & 0.6250 & 0.7000 \\
\bottomrule
\end{tabular}
\par\smallskip\begin{minipage}{\linewidth}\footnotesize These profiles retain the distinct contributions of pathway coverage, evidence support, and agreement with the requested answer.\end{minipage}
\end{table}


\paragraph{HA8: stable reactant reference and reactive conformer.}
Guided reports 29.46 kcal/mol and Standard 20.43 kcal/mol, against a reference of 29.48 kcal/mol. Their 9.03 kcal/mol difference comes from the reactant energy: Guided uses the lowest all-s-trans minimum, whereas Standard uses a reactive cisoid conformer. The IRC identifies the local conformer connected to the saddle; the task's barrier is measured from the specified stable reactant minimum. The two references can therefore be different conformers.

\paragraph{HH1: the same barrier with different IRC support.}
Both workflows report a binaphthyl racemization barrier of 23.48 kcal/mol; the reference is 23.01 kcal/mol. Their accuracy components are 0.973067 and 0.973064 and both graph components are 1. Guided has self-consistency 1, whereas Standard has 0.5. Their composites are consequently 0.9838 and 0.8338. Standard's archive contains candidate channels and IRC inputs; Guided completes the required chain from the saddle through the IRC to the endpoint calculations. This difference in path support contributes 0.1500 to the composite despite the matching barriers.

\paragraph{HI1: numerical agreement with different species and path coverage.}
Guided and Standard both have accuracy approximately 1. Their graph components are 0.999176 and 0.332509, and self-consistency components are 1 and 0.233333. The corresponding composites are 0.9999 and 0.7032. Matching numerical accuracy thus accompanies different coverage of the requested species, steps, and path calculations.

\paragraph{NU6: the optimized saddle and the intended second step.}
The Guided record gives $E_{\mathrm{TS}}=-247.030042392$ Ha for its concerted saddle region and $E_{\mathrm{betaine}}=-247.060060699$ Ha. Using 627.5094740631 kcal mol$^{-1}$ Ha$^{-1}$ gives
\begin{equation}
\left[E_{\mathrm{TS}}-E_{\mathrm{betaine}}\right]\,627.5094740631
=18.836772\ \mathrm{kcal\,mol^{-1}}.
\end{equation}
Standard reports 18.837 kcal/mol as the second-step ring-closure barrier. Guided leaves the betaine-to-product ring-closure saddle unresolved and describes the energy difference relative to the concerted saddle region. Its searches from interpolated betaine geometries and constrained-scan peaks repeatedly return toward the direct reactant-to-product saddle region. Assigning the energy difference to an elementary step therefore requires identifying the optimized saddle and its actual endpoints.

\subsection{Impossible-tier results and the requested negative conclusions}

Table~\ref{tab:si-imp-pairs} resolves Brief's mean composite of 0.482940 and Guided's 0.445060 into all ten task differences. Brief is higher on three tasks, Guided on five, and two tie. The Brief-minus-Guided mean is +0.037880, with the existing task-bootstrap interval [$-0.188220$, 0.283970]. Its weighted contributions are $-0.006607$ from graph, $-0.015500$ from self-consistency, and +0.060000 from accuracy. IA1 and IA9 provide the two large positive differences, while IG8 has the opposite accuracy ordering. Across the seven tasks with equal accuracy components in these two workflows, Brief-minus-Guided averages $-0.050000$.

\begin{table}[!htbp]
\centering
\small
\caption{Task-resolved comparison of Brief and Guided in the Impossible tier.}
\label{tab:si-imp-pairs}
\begin{tabular}{lrrrrr}
\toprule
Task & Guided C & Brief C & Brief $-$ Guided & Guided A & Brief A \\
\midrule
IA1 & 0.1717 & 1.0000 & +0.8283 & 0.0 & 1.0 \\
IA5 & 0.4000 & 0.3667 & -0.0333 & 0.0 & 0.0 \\
IA7 & 0.2583 & 0.2583 & +0.0000 & 0.0 & 0.0 \\
IA9 & 0.3667 & 0.9333 & +0.5666 & 0.0 & 1.0 \\
IB9 & 0.8200 & 0.7700 & -0.0500 & 0.7 & 0.7 \\
IC3 & 0.4000 & 0.4000 & +0.0000 & 0.0 & 0.0 \\
ID2 & 0.3339 & 0.1839 & -0.1500 & 0.0 & 0.0 \\
IE2 & 0.3333 & 0.1833 & -0.1500 & 0.0 & 0.0 \\
IF3 & 0.3667 & 0.4000 & +0.0333 & 0.0 & 0.0 \\
IG8 & 1.0000 & 0.3339 & -0.6661 & 1.0 & 0.0 \\
\bottomrule
\end{tabular}
\par\smallskip\begin{minipage}{\linewidth}\footnotesize C denotes the composite and A the accuracy component for the task-defined numerical or negative target. All values are retained from the publication score table.\end{minipage}
\end{table}


Brief reaches accuracy 1 for IA1, the square-geometry cyclobutadiene automerization target, and IA9, dioxetane thermolysis. Guided reaches accuracy 1 for IG8, the chromium-dimer target. Both have accuracy 0.7 in IB9, which contains positive and negative target components. Across all 40 Impossible runs, four have accuracy 1 and 14 have self-consistency 1 with accuracy 0. The following cases show why local calculation support and the conclusion about the requested state or pathway can score differently.

The preserved Brief/IC3 graph identifies the requested lowest open-shell singlet $^1A_2$ as inaccessible under the permitted route, marks its state-change edge accordingly, and withholds the requested singlet--triplet gap. It retains a separately characterized closed-shell $^1A_1$ branch. The PM7 value of +45.9 kcal/mol and GFN2-xTB value of $-1.7$ kcal/mol are labelled as exploratory relative energies for that branch, outside the requested gap. The mixed initial guess was used in a restricted calculation; the requested unrestricted electronic state remains untested. The final profile is Graph/SC/Accuracy = 1/1/0 and composite 0.4.

Record/IA5 distinguishes the requested disrotatory ring opening from an allowed conrotatory path. It marks the disrotatory edge inaccessible and reports the conrotatory path separately at 33.57 kcal/mol. The evaluation notes show that the matcher selected the positive path because it shared the endpoints, leaving accuracy at 0. Both the agent's stated conclusions and the scores used in the analysis are retained.

\subsection{Easy-tier energy references and delivery trajectories}
\label{sec:si-easy-case-audit}

The five Easy tasks examined here distinguish numerical reference choice, chemical identity matching, and the delivery of computed evidence. We inspected the retained runs' final mechanism graphs, native calculation logs, evaluation notes, and runtime events. Table~\ref{tab:si-easy-case-audit} reports the final scores. $\mathrm{SC}_{\mathrm{all}}$ grades the delivered graph, whereas $\mathrm{SC}_{\mathrm{Gold}}$ grades its reference-matched subgraph and excludes additional elements. Their difference identifies a change in the objects being assessed; the composite uses $\mathrm{SC}_{\mathrm{all}}$.

\begin{table}[!htbp]
\centering
\small
\setlength{\tabcolsep}{4pt}
\caption{Final component scores for the five inspected Easy tasks, rounded to four decimal places. Brief denotes Brief-errors. The two SC columns retain their separate assessment scopes. Accuracy is the final component after reference matching and the specified numerical penalty cap.}
\label{tab:si-easy-case-audit}
\begin{tabular}{llrrrrr}
\toprule
Task & Workflow & Graph & $\mathrm{SC}_{\mathrm{all}}$ & $\mathrm{SC}_{\mathrm{Gold}}$ & Accuracy & Composite \\
\midrule
F4 & Guided & 0.9167 & 1.0000 & 1.0000 & 1.0000 & 0.9917 \\
 & Standard & 0.5176 & 0.7500 & 0.0000 & 0.0000 & 0.2768 \\
 & Record & 0.6009 & 0.7500 & 0.0000 & 0.0000 & 0.2851 \\
 & Brief & 0.9167 & 1.0000 & 1.0000 & 1.0000 & 0.9917 \\
\addlinespace
G2 & Guided & 0.8889 & 0.5000 & 0.5000 & 0.0000 & 0.2389 \\
 & Standard & 1.0000 & 0.5000 & 0.5000 & 0.8782 & 0.7769 \\
 & Record & 1.0000 & 0.5000 & 0.5000 & 0.8782 & 0.7769 \\
 & Brief & 1.0000 & 0.5000 & 0.5000 & 1.0000 & 0.8500 \\
\addlinespace
D1 & Guided & 1.0000 & 1.0000 & 1.0000 & 0.7000 & 0.8200 \\
 & Standard & 0.7917 & 0.5000 & 0.5000 & 1.0000 & 0.8291 \\
 & Record & 0.7917 & 1.0000 & 1.0000 & 0.7000 & 0.7992 \\
 & Brief & 1.0000 & 0.7500 & 0.7500 & 0.7000 & 0.7450 \\
\addlinespace
D3 & Guided & 1.0000 & 0.7500 & 0.7500 & 0.7000 & 0.7450 \\
 & Standard & 1.0000 & 0.7500 & 0.7500 & 0.9995 & 0.9247 \\
 & Record & 0.9444 & 0.3500 & 0.3500 & 0.9909 & 0.7940 \\
 & Brief & 0.8333 & 0.7500 & 0.7500 & 0.7000 & 0.7283 \\
\addlinespace
L2 & Guided & 1.0000 & 0.7500 & 0.7500 & 0.7000 & 0.7450 \\
 & Standard & 1.0000 & 0.5000 & 0.5000 & 0.7000 & 0.6700 \\
 & Record & 1.0000 & 0.7500 & 0.7500 & 0.8832 & 0.8549 \\
 & Brief & 1.0000 & 0.8333 & 0.7500 & 0.8832 & 0.8799 \\
\bottomrule
\end{tabular}
\end{table}

\paragraph{F4: adduct identity in the final mechanism graph.}
Guided delivered hydride transfer from the Lewis adduct of borane and formaldehyde to methoxyborane. Its bound reactant and saddle energies were $-141.041124115$ and $-141.015601446$~hartree, giving 16.01572~kcal~mol$^{-1}$ and supporting the declared 16.02. Independent saddle frequency analysis gave one imaginary frequency, $-398.7218$~cm$^{-1}$. Both actual IRC outer frames subsequently converged to the corresponding zero-imaginary-frequency endpoint minima. The four workflows obtained essentially the same saddle and adduct energies.

The F4 reference defines hydride transfer from the bonded Lewis adduct, represented as \texttt{[BH3-][O+]=C}, with a 16.02-kcal~mol$^{-1}$ barrier. Guided and Brief declared this adduct in their final graphs and obtained Accuracy close to 1. Standard and Record retained the separated-fragment representation \texttt{[BH3].C=O} for the reactant node. Their final graph declarations therefore did not identify the reference adduct and received Accuracy~0, although their native calculations supported essentially the same physical pathway. The comparison separates the chemical identity declared in the delivered graph from the path established by the calculations.

\paragraph{G2: computed evidence and final structured delivery.}
Guided's native logs contain a concerted double-H-transfer saddle at $-189.093172517$~hartree with one imaginary frequency, $-444.9504$~cm$^{-1}$. The optimized cis reactant complex and the IRC-derived reactant minimum have energies of $-189.094518484$ and $-189.094518735$~hartree. They give barriers of 0.84461 and 0.84476~kcal~mol$^{-1}$, respectively, close to the 0.87 reference.

The runtime recorded four rejections of a numerical barrier on a companion complex-product edge whose scientific checks remained open. Meanwhile, delivery reports for the selected separated-products target returned \texttt{complete=true} with no unresolved target. After an initial deletion, a Supervisor request prompted restoration of all numerical barriers and the saddle's energy/log fields. The fourth barrier rejection was followed by deletion of the companion-edge values. A subsequent artifact check rejected the saddle-evidence claim; the next recorded edit removed the selected edge's remaining numerical barrier and the saddle's energy/log fields. The successful write and program output confirm these removals. The final graph retained the values in descriptive text, so the numerical target received Accuracy~0. Its graded SC contribution remained 0.5.

Brief retained a numerical barrier of 0.844 and obtained Accuracy~0.99999. Standard and Record declared approximately 3.4 using a lower-energy reactant reference; Record's native report identifies this as a trans-diimide stack and separately records the approximately 0.84 cis-reference result.

\paragraph{D1 and D3: numerical penalties with different physical origins.}
The evaluation notes explicitly attribute Accuracy~0.7 to the maximum numerical-deviation penalty of 0.3. For D1, 15.738 versus 18.830~kcal~mol$^{-1}$ gives an underlying Gaussian score of 0.5880 before this cap. Guided and Standard obtained nearly identical saddle energies, but their reactant energies differ by 3.11859~kcal~mol$^{-1}$. Their butadiene dihedrals are approximately $26.33^\circ$ and $-179.93^\circ$, corresponding to reactive s-cis and lower s-trans references. Record also delivered the s-cis barrier and retained SC~1 with Accuracy~0.7.

D3 has a different origin. Guided and Standard delivered barriers of 18.14687 and 14.55472~kcal~mol$^{-1}$. Their reactant energies differ by only 0.09899~kcal~mol$^{-1}$, whereas their saddle energies differ by 3.69113~kcal~mol$^{-1}$. Both saddle logs contain exactly one imaginary frequency. The archived evaluations separately assign SC~0.75 with product-endpoint association warnings; the Accuracy cap follows the numerical barrier comparison.

\paragraph{L2: two reference barriers and the selected comparison edge.}
Guided used the reactive cZc minimum and obtained 18.98740~kcal~mol$^{-1}$; its graph states that the stable-to-reactive conformer contribution was uncomputed. Record and Brief used the lower tZt reference, giving approximately 28.275. The conformer energy difference is 9.28761~kcal~mol$^{-1}$.

Standard delivered both references and both barriers: 18.99 on edge \texttt{e1} and 28.28 on \texttt{e2}. The archived evaluation selected \texttt{e1}. Both candidates have identical bond-change lists, and the unspecified-stereochemistry reference gives equal endpoint substitution costs for these reactants. The final edge-identifier tie-break selects \texttt{e1}, so the local-reference value receives the capped Accuracy~0.7. Standard's final report explicitly contains the stable-reference barrier and conformer gap.

\subsection{Sensitivity to target-object assignment in Impossible tasks}
\label{sec:matcher-sensitivity}

Record/IA5 and Brief/IC3 distinguish pathways or electronic states that share the same endpoint connectivity. Using their original final graphs and process records, we calculated how the Impossible-tier means would change if the existing negative-credit rule were applied to the requested pathway or state. The 200-run evaluation table remains the reported dataset; these conditional results are listed separately. Graph and Self-consistency remain fixed in both cases, as do all scores for the other 198 runs.

IA5 requests the thermal disrotatory ring opening of cyclobutene. Record's final graph marks the disrotatory edge \texttt{e\_dis} and its transition-state node inaccessible, leaving its barrier unspecified. A separate conrotatory edge, \texttt{e\_con}, carries a barrier of 33.57 kcal/mol. The final process record explicitly distinguishes the requested disrotatory target from this computed alternative. Both edges have the same connectivity endpoints and bond-change signature; identifying the requested pathway also requires the conrotatory/disrotatory distinction.

IC3 requests the lowest singlet state of phenylnitrene and its singlet--triplet gap. Brief's original graph distinguishes the closed-shell singlet \texttt{S1\_1A1} from the open-shell target \texttt{S2\_1A2}. It marks the latter node and its state-change edge inaccessible, withholds the requested numerical gap, and explicitly excludes the closed-shell branch's relative energies from that gap. The two state-change edges again have identical endpoint SMILES and bond-change signatures. The final baseline for Brief/IC3 has Graph $=1$, Self-consistency $=1$, Accuracy $=0$, and composite $=0.4$.

In both graphs, the matcher selects the positive alternative pathway or state. Matching the requested negative target to its reference would earn credit under the unchanged negative-claim rule. A correction must first use the run's final statements to distinguish the requested pathway or state from the positive alternative. It must also retain any contradictory positive claim about that same target.

Both cases in the final evaluation have Graph $=1$, Self-consistency $=1$, Accuracy $=0$, and composite $=0.4$. Assigning Accuracy $=1$ to either intended negative target changes that case's conditional composite by $0.6$, following the fixed Accuracy weight. Its configuration's ten-task Impossible mean therefore increases by $0.06000$. Table~\ref{tab:matcher-sensitivity} reports the four scenarios; the accompanying \texttt{matcher\_sensitivity\_scenarios.csv} retains their configuration-level means and ranks.

\begin{table}[htbp]
\centering
\small
\caption{Impossible-tier mean composite scores under two conditional target-assignment corrections. Only the indicated Accuracy values change; Graph, Self-consistency, and the other 198 runs remain fixed. The baseline row gives the final evaluation values.}
\label{tab:matcher-sensitivity}
\begin{tabular}{@{}lrrrr@{}}
\toprule
Scenario & Guided & Standard & Record & Brief \\
\midrule
Final evaluation values & 0.44506 & 0.24733 & 0.23250 & 0.48294 \\
Record/IA5 target credit & 0.44506 & 0.24733 & 0.29250 & 0.48294 \\
Brief/IC3 target credit & 0.44506 & 0.24733 & 0.23250 & 0.54294 \\
Both target credits & 0.44506 & 0.24733 & 0.29250 & 0.54294 \\
\bottomrule
\end{tabular}
\end{table}

The Impossible-tier ordering Brief $>$ Guided $>$ Record remains the same in all scenarios. Crediting Record's IA5 target moves Record above Standard; crediting Brief's IC3 target leaves Brief first. These assignments can therefore reverse the Record--Standard ordering while preserving every Guided--Standard score, paired difference, and primary comparison. The calculations apply the two stated score substitutions; the final evaluation table and scientific-validation results remain frozen.


\subsection{Discussion of the supplementary comparisons}

Tool and task graphs organize computational execution. Cloudra's Delivery DAG constrains the agents' submitted mechanism representation through species-occurrence nodes and elementary-step edges that carry molecular geometries, bonding patterns, and reaction connectivity. Cheminformatics algorithms establish and compare identities and atom correspondences; frequency and IRC/endpoint calculations determine the supported stationary points and connections. The Normal and Hard cases show why these assignments matter: HA8 requires correcting the reactant energy reference, and NU6 requires recognizing that a search has returned to a different saddle region. Agents use the represented species and steps to choose new geometries, continue a path calculation, or revise the mechanism, complementing external planning and tool feedback.\citep{valmeekam2023planning,ref18,huang2024selfcorrect}

P2 retains one run per task and configuration. Its task-cluster bootstrap intervals measure variation across the sampled tasks. P1b supplies three repeats on each of four tasks under each configuration; these 48 runs provide repeat measurements on those four tasks within the 50-task P2 study. The P2 table includes all 200 runs and their three components. Some Hard tasks score qualitative channel identity, so the gain includes pathway identification and completion of the supporting calculations as well as numerical barriers. Repeating the full task set would measure run-to-run variation and how often the agents select each route.

The Impossible rubric concerns targets judged physically nonexistent under the specified permitted methods; self-consistency credit records the supported negative conclusion. A positive step requires its own TS--IRC--endpoint chain. IA5 and IC3 require matching the requested pathway and electronic state before scoring the answer. Brief's crossed outcomes suggest task-dependent effects of diagnostic content; the causal effect of diagnostic length remains unmeasured. In NU6, the matching energy difference prompts a check of the assigned step. The native geometries and paths determine whether two calculations found the same saddle.

Timing comparisons use batch elapsed time for P1b and accumulated job-hours for P2. Separate measurements of queueing, calculation, agent-active time, and end-to-end P2 latency remain future work.


Cloudra's default preorganization samples 32 rigid relative placements, relaxes them with MMFF94, and selects candidates by energy and clustering. It samples local arrangements; global-minimum coverage requires a separate assessment. CREST is an optional backend, and its coverage across the 200 runs remains to be established. Flexible systems may need broader sampling and more quantum-chemical calculations. Further work should measure which conformers are found and how calculations can be allocated among competing candidates.

The benchmark scores electronic energy differences alone; zero-point, thermal/entropic, standard-state, and solvation free-energy contributions remain outside this comparison. NI5 specifies weakly acidic aqueous solution, and NU5 specifies dichloromethane at high concentration; both are scored against the prescribed electronic-energy references. Solution-phase free-energy barriers and solvent sensitivity, including the magnitude and direction of solvent effects, remain unassessed.

All workflow comparisons used the same atom-mapping coefficients and ambiguity window, and the same branch-crowding decay, depth, and congestion settings. Their individual effects remain unmeasured. The imaginary-mode alignment ratio of 0.5 is advisory; physical validation uses the stationary-point and IRC/endpoint calculations. All P2 configurations used DeepSeek-v4.1-flash (1017), the API alias \texttt{deepseek-v4-flash}, and the recorded request settings. Additional backbones would test whether the workflow improvements extend to other models.


The Easy cases show two different problems. In F4, the final graphs assign different reactant identities despite calculations supporting essentially the same pathway. In G2, validation feedback prompts deletion of structured numerical fields while the values remain in prose. Recording validation status separately would let the agent state a quantity's scientific qualification while retaining the quantity and its source. The conditional IA5/IC3 analysis changes which pathway or electronic state is scored while keeping the validation rules and final publication dataset fixed.
